\clearpage
\SourceLeaf{587}\pdfbookmark[0]{Reviews of Works by Other Authors.}{anzeigen-nicht-eigner}
\thispagestyle{empty}
\vspace*{35mm}
\begin{center}
{\LARGE REVIEWS}\par\vspace{16mm}
{\large OF WORKS BY}\par\vspace{16mm}
{\LARGE OTHER AUTHORS.}
\end{center}
\clearpage
\SourceLeaf{589}\fancyhead[L]{\small REVIEWS OF WORKS BY OTHER AUTHORS.}
\pdfbookmark[1]{On Dalton’s Theory. 1830.}{daltonsche-theorie}
\begin{center}\small Göttingische gelehrte Anzeigen.\quad 1830 December 11.\end{center}

\textit{On \textup{\textsc{Dalton}}’s Theory, by \textup{J. F.\ \textsc{Benzenberg}.} Düsseldorf, \textup{1830.} Published by J. E.\ Schaub.} (192 pages in 8vo, with three lithographed plates).

Few physicists have so far accepted the hypothesis put forward by \textsc{Dalton}, that the different gases composing atmospheric air exert no pressure at all on one another but constitute as many atmospheres, as it were independent of one another. Among those who have accepted it, however, Mr.\ \textsc{Benzenberg} stands out particularly for the tireless and heartfelt zeal with which he has defended that hypothesis for nearly twenty years. In particular, he found an important argument \textit{for} \textsc{Dalton}’s hypothesis in \textsc{Daubuisson}’s trigonometric and barometric measurement of the height of Monte Gregorio. Clearly, if \textsc{Dalton}’s hypothesis is true, barometric measurements of heights must be calculated differently than under the usual theory. For Monte Gregorio, which is 5260 feet high, Mr.\ \textsc{Benzenberg} found that the first calculation gave a result 16 feet smaller than the other one. The latter result exceeded that of the trigonometric measurement by very nearly the same amount. Thus, according to Mr.\ \textsc{Benzenberg}’s calculation, adopting \textsc{Dalton}’s hypothesis would remove the difference almost entirely. Mr.\ \textsc{Benzenberg} first published these calculations in \textsc{Gilbert}’s Annalen der Physik in 1812 and subsequently
\SourceLeaf{590}presented them repeatedly elsewhere. Mr.\ \textsc{Benzenberg} also discussed other branches of physics connected with \textsc{Dalton}’s way of thinking, such as acoustics and eudiometry. His aim was less to seek arguments \textit{for} the hypothesis there than to combat the arguments that can be, and have been, drawn from them against it. All of this is presented together again in the present work.

The recognition. that these efforts have received from physicists and mathematicians seems to have fallen short, at least where those abroad are concerned. Mr.\ \textsc{Benzenberg} relates in his preface that during his stay in Paris in 1815 he could find no copy of \textsc{Gilbert}’s Annalen at the library of the Institute, and that the copy at the great Royal Library was uncut. He says that he spoke with \textsc{Laplace} about his theory for calculating barometer heights under \textsc{Dalton}’s hypothesis, but that this great geometer was already old at the time.

German physicists, however, cannot fairly be accused of having ignored the subject. In the article Atmosphere in the new edition of the dictionary of physics, we read that ‘Mr.\ \textsc{Benzenberg} is the most distinguished, thorough and zealous defender of \textsc{Dalton}’s theory, that heights are calculated more accurately with its application than without it, and that this provides strong evidence that the theory is correct.’

But the author of this article seems never to have \textit{checked} \textsc{Benzenberg}’s calculations. All the numbers were simply copied from \textsc{Gilbert}’s Annalen. The same applies to what is said of those calculations in the article Measurement of Heights in the dictionary just mentioned. Perhaps the authors of both articles regarded checking as less essential because, as they placed less trust than Mr.\ \textsc{Benzenberg} in barometric measurements of heights generally, they did not in the end acknowledge the force of his calculations as evidence. The author of the first article, to be sure, thereby partly retracts what was quoted above.

The present review cannot aim to develop our own view of whether \textsc{Dalton}’s hypothesis itself is tenable. Instead we confine ourselves to what Mr.\
\SourceLeaf{591}\textsc{Benzenberg} has put forward again in the present work in support of it, and in particular to his calculation of barometric measurements of heights.

First of all, it seemed necessary to check whether \textsc{Benzenberg}’s calculation itself was correct. To our surprise, it turned out that the calculation is wrong, and that a correctly performed calculation gives exactly the opposite result.

We shall have to demonstrate this in some detail here and far enough to enable everyone to judge for themselves, all the more because the issue is less an error of calculation in the strict sense than an error in reasoning.

On the second and third pages of the book we find two tables, each of three columns, which we must examine one by one. (The same tables appear in \textsc{Gilbert}’s Annalen der Physik, vol.\ 42, pp.\ 163 and 164.)

The first column gives, according to \textsc{Dalton}, the amounts in 100 parts of dry atmospheric air:
\begin{center}
\begin{tabular}{lclrl}
Nitrogen & .\quad .\quad . & as & $78.93$ & parts\\
Oxygen & . & — & $21.00$ & —\\
Carbon dioxide & & — & $0.07$ & —
\end{tabular}
\end{center}
all by volume. (The book itself has $21.90$ in place of the second number, which is an obvious printing error.)

The second column gives the specific gravities
\begin{center}
\begin{tabular}{lcr}
of nitrogen & .\quad .\quad . & $0.9691$\\
of oxygen & .\quad . & $1.1148$\\
of carbon dioxide & & $1.5000$
\end{tabular}
\end{center}
with that of ordinary atmospheric air taken as unity. The first and third numbers are taken from \textsc{Biot}. As he says himself, Mr.\ \textsc{Benzenberg} calculated the second ‘so that the mixture has exactly the weight of dry atmospheric air given by \textsc{Biot}, namely $\frac{1}{10493}$ of that of mercury at $0^0$ Réaumur and at a pressure of 28 inches.’ We can see that Mr.\ \textsc{Benzenberg} had not made clear to himself what he actually meant to calculate here, since what he has just stated is extraneous to \textit{this} calculation. Clearly the sole aim was to give
\SourceLeaf{592}the mixture a specific gravity of 1. In any case, this lack of clarity has no further effect here. The calculation, however, was not performed very accurately. The stated data give a specific gravity for oxygen gas different from $1.1148$, namely $1.11447$. This error. is, however, quite insignificant.

The third column gives the parts by \textit{weight} of the individual gases contained in 100 parts of dry air, namely
\begin{center}\begin{tabular}{lcrl}
Nitrogen & .\quad .\quad . & $76.49$ & parts\\
Oxygen & . & $23.41$ & —\\
Carbon dioxide & & $0.10$ & —
\end{tabular}\end{center}
These numbers are obviously just the products of the numbers in the first column and the corresponding numbers in the second.

The fourth column gives the height at which each individual atmosphere would hold the barometer (assuming \textsc{Dalton}’s hypothesis is correct), or the share of the total pressure, with the latter taken as a mercury column of $27.76$ inches for dry air. Mr.\ \textsc{Benzenberg}’s numbers are
\begin{center}\begin{tabular}{llcrl}
for the & nitrogen atmosphere & .\quad .\quad . & $21.2326$ & inches\\
—\quad — & oxygen atmosphere & . & $6.4986$ & —\\
—\quad — & carbon dioxide atmosphere & & $0.0278$ & —
\end{tabular}\end{center}
We shall return to the calculation of this column in a moment.

The fifth column gives the specific gravities of the gases relative to mercury at the freezing point and at a pressure of 28 inches of mercury. These numbers, namely $\frac{1}{10830}$, $\frac{1}{9414}$ and $\frac{1}{6997}$ respectively, are simply the products of the numbers in the second column and $\frac{1}{10495}$.

The sixth column is headed only \textit{Constant Number.} But we can see that what is meant is the so-called subtangent, or the height a hypothetical atmosphere of uniform density equal to that of the actual atmosphere at the bottom would have to reach to exert the same pressure. For ordinary dry air this number is therefore the product of 10495 and 28 inches, or $24488\frac{1}{3}$ feet. For the three separate atmospheres as \textsc{Dalton} conceives them, the numbers will likewise be the products of 28 inches and the denominators of the fractions in the fifth column. More simply, they can be found by dividing $24488.33$ feet by the numbers in the second column. We set out these numbers here,
\SourceLeaf{593}both as Mr.\ \textsc{Benzenberg} gives them and as they follow from a more precise calculation
\begin{center}\begin{tabular}{lcrlrl}
& & \multicolumn{2}{c}{from Mr.\ \textsc{Benz}.} & \multicolumn{2}{c}{from more precise calc.}\\
Nitrogen & .\quad . & $25270$ & feet & $25269.15$ & feet\\
Oxygen & . & $21966$ & — & $21973.01$ & —\\
Carbon dioxide & & $16326$ & — & $16325.56$ & —
\end{tabular}\end{center}

None of the objections raised so far against Mr.\ \textsc{Benzenberg} is of any consequence. We now come to the essential point, the way the fourth column was calculated. Mr.\ \textsc{Benzenberg} gives no explanation at all of \textit{how} he made the calculation. He merely says that these are examples of partnership arithmetic. It is easy to see, however, that he simply made the numbers in the fourth column proportional to those in the third, or derived the former from the product of $27.76$ inches and the following numbers.
\begin{center}\begin{tabular}{rl}
$0.7649$ & for nitrogen\\
$0.2341$ & for oxygen\\
$0.0010$ & for carbon dioxide
\end{tabular}\end{center}

And this is incorrect.

For in \textsc{Dalton}’s hypothesis, the total pressure of the nitrogen atmosphere relative to that of the oxygen atmosphere follows the \textit{compound} ratio of the shares by weight that the two gases have in the lowest cubic foot of mixed air and of their respective subtangents. The shares by weight alone would give the wrong ratio. The pressures must therefore be proportional to the \textit{products} of the numbers in the third and sixth columns, or equivalently to the quotients obtained by dividing the numbers in the third column by those in the second, hence simply to the numbers in the first column.

A little reflection makes this clear in itself, for the meaning of the numbers in the first column can also be expressed as follows. The nitrogen contained in a volume of 100 parts of ordinary dry air at the bottom of the atmosphere would, taken on its own at the same pressure as that air, occupy a space of only $76.93$ parts, and the same
\SourceLeaf{594}holds for the other gases. Thus, with each of these three gases now spread throughout the space of 100 parts and regarded as independent of the others, the mercury pressures that they each balance are in the following ratio
\begin{center}\begin{tabular}{rl}
$78.93$ & for nitrogen\\
$21.00$ & for oxygen\\
$0.07$ & for carbon dioxide.
\end{tabular}\end{center}

We can see, then, that in his fourth column Mr.\ \textsc{Benzenberg} apportioned the total pressure of $27.76$ inches incorrectly. The correct numbers are
\begin{center}\begin{tabular}{rlll}
$21.9110$ & inches & for & nitrogen\\
$5.8296$ & — & — & oxygen\\
$0.0194$ & — & — & carbon dioxide.
\end{tabular}\end{center}
The numbers in the fourth and sixth columns are the quantities from which the pressure of the atmosphere at any height must be calculated under \textsc{Dalton}’s hypothesis. Naturally, the corrected values therefore give different results.

\Needspace{10\baselineskip}
We have taken the trouble to carry out this calculation for several heights. The following table gives the results.
\begin{center}\begin{tabular}{r|r|r|r}
& \multicolumn{3}{c}{Barometer height in inches}\\
\multicolumn{1}{c|}{Height} & \multicolumn{1}{c|}{under usual} & \multicolumn{2}{c}{under \textsc{Dalton}’s hypothesis}\\
\multicolumn{1}{c|}{feet} & \multicolumn{1}{c|}{theory} & \multicolumn{1}{c|}{our calc.} & \multicolumn{1}{c}{Mr.\ \textsc{Benz}.’s calc.}\\\hline
$5000$ & $22.6332$ & $22.6350$ & $22.6179$\\
$10000$ & $18.4532$ & $18.4589$ & $18.4314$\\
$15000$ & $15.0452$ & $15.0555$ & $15.0221$\\
$20000$ & $12.2666$ & $12.2814$ & $12.2458$
\end{tabular}\end{center}

For the barometer heights under the usual theory we have used the numbers from our own calculation here. The values Mr.\ \textsc{Benzenberg} calculated for the same heights, also under the usual theory (p.\ 12), differ somewhat in some cases. We can vouch for the correctness of our calculation.

Thus the barometer heights under \textsc{Dalton}’s hypothesis differ from those under the usual theory by the following amounts.
\SourceLeaf{595}\begin{center}\begin{tabular}{rl|rl|rl}
& & \multicolumn{4}{c}{Difference in barometer height}\\
\multicolumn{2}{c|}{Height} & \multicolumn{2}{c|}{our calc.} & \multicolumn{2}{c}{Mr.\ \textsc{Benz}.’s calc.}\\\hline
$5000$ & feet & $+0.0018$ & inches & $-0.0153$ & inches\\
$10000$ & — & $+0.0057$ & — & $-0.0218$ & —\\
$15000$ & — & $+0.0103$ & — & $-0.0231$ & —\\
$20000$ & — & $+0.0148$ & — & $-0.0208$ & —
\end{tabular}\end{center}

At a given height, \textsc{Dalton}’s hypothesis therefore gives a higher barometer reading than the usual theory, contrary to Mr.\ \textsc{Benzenberg}’s view that it gives a lower one. Correspondingly, a given barometer reading yields a greater height under that hypothesis, whereas his conclusion gives a smaller one. For Monte Gregorio he gives a difference of $-16$ feet. The correct difference is $+2$ feet. At lower heights the difference is very nearly proportional to the square of the height. A straightforward consideration could easily have predicted this \textit{a priori}, but for brevity we shall leave it undeveloped here. Mr.\ \textsc{Benzenberg}’s differences, by contrast, are nearly proportional to the heights themselves when these are small. That alone was enough to show that they were incorrect. The numerical calculation above, moreover, was carried out for dry air only. Taking water vapour into account would change the relative differences very little.

We add a few remarks to this account.

I.\quad Our result is that the difference between the barometer height under \textsc{Dalton}’s hypothesis and that calculated in the usual way is \textit{positive} and, at moderate heights, becomes nearly proportional to their squares. This is a general result and is independent of the assumed specific gravities of the individual gases that make up the mixed air. It would therefore be futile to expect a more favourable result from other values of these specific gravities.

II.\quad As early as 1807, \textsc{Tralles} gave a correct account of how to calculate barometer heights under \textsc{Dalton}’s hypothesis. Even a cursory look shows that his result agrees fully with ours in essentials. It is therefore surprising that the author of one of the articles in the dictionary of physics mentioned above claims that \textsc{Tralles} finds very nearly the same difference as Mr.\ \textsc{Benzenberg}. In fact, it differs both in sign and in the law it follows. At moderate
\SourceLeaf{596}heights the proportionality is to the square of the height. It does not hold for the height itself. Perhaps this oversight is due to a misprint, though one easily recognised as such. In \textsc{Tralles}’s paper (\textsc{Gilbert}’s Annalen, vol.\ 27, p.\ 445), $\frac{B}{b+\beta}$ appears once in place of $\frac{b+\beta}{B}$. Yet the fraction is given correctly again where it is applied to a specific example (same page, bottom line), and the difference is entered with the correct sign on the following page. Still more surprising is that Mr.\ \textsc{Benzenberg} failed to compare the result of his predecessor, known as a thorough mathematician, with his own, to which it is entirely opposed. One would indeed suppose that Mr.\ \textsc{Benzenberg} did not know the paper at all, were it not for his explicit mention of it, though he merely dismisses it on p.\ 15. ‘Mr.\ \textsc{Tralles} used algebraic calculation. This is unnecessary. If you picture four barometers, you can do it by the rule of three, and you need no learning at all’. This principle, which Mr.\ \textsc{Benzenberg} has loudly professed on many occasions — we shall leave aside whether they were always the right ones — may be granted in the present case, and our account itself may then serve to confirm it, if we have managed to make it sufficiently clear.

III.\quad Trying to remove the discrepancy between the barometric and trigonometric measurements of Monte Gregorio by means of \textsc{Dalton}’s hypothesis is futile (under it the discrepancy even increases by two feet). It is thus an established fact that one or both measurements lack the accuracy Mr.\ \textsc{Benzenberg} thought he could ascribe to them. In our view, all three possible sources of error may have played a part. First, the fluctuations in the measured barometer heights themselves. Second, the constant used in the calculation, which Mr.\ \textsc{Benzenberg} had based on \textsc{Biot}’s weighing of atmospheric air. It can probably be determined much more reliably by making appropriate use of numerous mountain heights measured both barometrically and trigonometrically. Third, the trigonometric measurement of Monte Gregorio itself may also have contributed to the discrepancy, since we cannot fully agree with Mr.\ \textsc{Benzenberg}’s judgment that its accuracy is unsurpassable. The determination of the dist\SourceLeaf{597}ance from the summit to the station at the foot of the mountain was based on only a short baseline. Its length (670 metres) does seem to have been measured with great care (four days were spent on it), but the angle opposite it (only $6^0\;14'$) was obtained solely by inference from the other two angles, with no measurement of its own. Such a procedure always requires great caution even when excellent instruments are used. Yet \textsc{Daubuisson}’s angle-measuring instrument, a \textsc{Lenoir} repeating circle eight inches in diameter, seems to have been a very mediocre one. We see that among the ten repetitions that make up the angle measurements at each of the two stations, the individual pairs on one occasion give differences of more than a minute for the single angle. Nor does \textsc{Daubuisson} say anything about the nature of the signals he used as targets, so we cannot judge how precisely they could be sighted for the intersection or whether there might have been an unfavourable phase of illumination. An error of half a minute in the inferred third angle would already change the measured height by seven feet. — We may leave the matter there, for the explanation of the discrepancy is quite immaterial to the present question. It is enough that \textsc{Dalton}’s hypothesis can contribute nothing to it.

